\documentclass[11pt,a4paper]{article}
% Preâmbulo comum. Não duplicar pacotes nos módulos.
\usepackage[utf8]{inputenc}
\usepackage[T1]{fontenc}
\usepackage[brazil]{babel}
\usepackage{amsmath,amssymb,amsthm,mathtools}
\usepackage{geometry}
\geometry{a4paper,margin=2.3cm}
\usepackage{microtype}
\usepackage{enumitem}
\usepackage{graphicx}
\usepackage{xcolor}
\usepackage{booktabs}
\usepackage{hyperref}
\hypersetup{colorlinks=true,linkcolor=blue!40!black,urlcolor=blue!50!black}

\setcounter{tocdepth}{2}
\numberwithin{equation}{section}

\theoremstyle{plain}
\newtheorem{teorema}{Teorema}[section]
\newtheorem{proposicao}[teorema]{Proposição}
\newtheorem{lema}[teorema]{Lema}
\newtheorem{corolario}[teorema]{Corolário}

\theoremstyle{definition}
\newtheorem{definicao}[teorema]{Definição}
\newtheorem{exemplo}[teorema]{Exemplo}
\newtheorem{exercicio}{Exercício}[section]

\theoremstyle{remark}
\newtheorem{observacao}[teorema]{Observação}
\newtheorem{contraexemplo}[teorema]{Contraexemplo}
\newtheorem{intuicao}[teorema]{Intuição}

\newcommand{\R}{\mathbb{R}}
\newcommand{\N}{\mathbb{N}}
\newcommand{\Z}{\mathbb{Z}}
\newcommand{\Q}{\mathbb{Q}}
\DeclareRobustCommand{\eps}{\varepsilon}
\DeclareRobustCommand{\vet}[1]{\vec{#1}}
\DeclareRobustCommand{\dif}{\mathrm{d}}
\DeclareRobustCommand{\abs}[1]{\left\lvert #1 \right\rvert}
\DeclareRobustCommand{\norm}[1]{\left\lVert #1 \right\rVert}
\DeclareMathOperator{\grad}{grad}
\DeclareMathOperator{\Div}{div}
\DeclareMathOperator{\rot}{rot}

\newcommand{\base}[1]{\par\medskip\noindent\textbf{Base para relembrar.} #1\par\medskip}
\newcommand{\cuidado}[1]{\par\medskip\noindent\textbf{Quando a hipótese falha.} #1\par\medskip}

\newcommand{\figura}[2]{%
  \begin{figure}[htbp]
  \centering
  \IfFileExists{#1}{%
    \includegraphics[width=0.84\textwidth]{#1}%
  }{%
    \fbox{\parbox{0.8\textwidth}{\centering Figura ainda não gerada:\\ \texttt{#1}}}%
  }
  \caption{#2}
  \end{figure}%
}

\newenvironment{esqueleto}{%
  \par\medskip\noindent\textbf{Esqueleto da demonstração.}\par\nopagebreak
  \begin{enumerate}[leftmargin=1.4em]
}{\end{enumerate}\medskip}

\title{Solu\c{c}\~ao do caderno --- $\eps$-$\delta$}
\author{}
\date{}
\begin{document}
\maketitle
\tableofcontents

As seis provas s\~ao independentes do texto das notas. Cada uma come\c{c}a da defini\c{c}\~ao.

\section{Limite linear}

\textbf{Enunciado.} $\lim_{x\to 2}(3x-1)=5$, com $\delta$ expl\'icito.

\begin{proof}
$\abs{(3x-1)-5}=\abs{3x-6}=3\abs{x-2}$. Dado $\eps>0$, escolha $\delta=\eps/3$. Se $0<\abs{x-2}<\delta$, ent\~ao
\[
\abs{(3x-1)-5}=3\abs{x-2}<3\cdot\frac{\eps}{3}=\eps.
\]
\end{proof}

\section{Limite do quadrado}

\textbf{Enunciado.} $\lim_{x\to 3}x^2=9$, com $\delta$ expl\'icito.

\begin{proof}
$\abs{x^2-9}=\abs{x-3}\abs{x+3}$. Se $\abs{x-3}<1$, ent\~ao
\[
\abs{x+3}=\abs{(x-3)+6}\le \abs{x-3}+6<7.
\]
Seja $M=7$. Dado $\eps>0$, escolha $\delta=\min(1,\eps/7)$. Se $0<\abs{x-3}<\delta$, valem $\abs{x-3}<1$ e $\abs{x-3}<\eps/7$, logo
\[
\abs{x^2-9}=\abs{x-3}\abs{x+3}<7\abs{x-3}<7\cdot\frac{\eps}{7}=\eps.
\]
\end{proof}

\section{Unicidade}

\textbf{Enunciado.} Se os dois limites finitos existem, eles coincidem.

\begin{proof}
Seja $r>0$ tal que $f$ est\'a definida em $(a-r,a+r)\setminus\{a\}$. Suponha $L\neq M$ e ponha $\eps=\abs{L-M}/2>0$. Existem $\delta_1>0$ e $\delta_2>0$ tais que
\begin{align*}
0<\abs{x-a}<\delta_1 &\Rightarrow \abs{f(x)-L}<\eps,\\
0<\abs{x-a}<\delta_2 &\Rightarrow \abs{f(x)-M}<\eps.
\end{align*}
Seja $\delta=\min(\delta_1,\delta_2,r)$ e escolha $x=a+\delta/2$. Ent\~ao $0<\abs{x-a}<\delta$, o ponto est\'a no dom\'inio, e
\begin{align*}
\abs{L-M}
&\le \abs{f(x)-L}+\abs{f(x)-M}
<\eps+\eps
=\abs{L-M}.
\end{align*}
A desigualdade estrita $\abs{L-M}<\abs{L-M}$ \'e falsa. Logo $L=M$.
\end{proof}

\section{Confronto}

\textbf{Enunciado.} Se $f\le g\le h$ numa vizinhan\c{c}a furada e os extremos tendem a $L$, ent\~ao $g$ tende a $L$.

\begin{proof}
Dado $\eps>0$, existe $\delta_1>0$ tal que $0<\abs{x-a}<\delta_1$ implica $\abs{f(x)-L}<\eps$, isto \'e, $L-\eps<f(x)<L+\eps$. Existe $\delta_2>0$ tal que $0<\abs{x-a}<\delta_2$ implica $L-\eps<h(x)<L+\eps$. Seja $\delta=\min(\delta_0,\delta_1,\delta_2)$. Se $0<\abs{x-a}<\delta$, ent\~ao
\[
L-\eps<f(x)\le g(x)\le h(x)<L+\eps,
\]
logo $\abs{g(x)-L}<\eps$.
\end{proof}

\section{Deriv\'avel implica cont\'inua}

\textbf{Enunciado.} Se $f'(a)$ existe, ent\~ao $\lim_{x\to a}f(x)=f(a)$.

\begin{proof}
Por defini\c{c}\~ao, $f'(a)=\lim_{h\to 0}(f(a+h)-f(a))/h$, com $h\neq 0$. Para esses $h$, escreva
\[
f(a+h)-f(a)=q(h)\,h, \qquad q(h)=\frac{f(a+h)-f(a)}{h}.
\]
O alvo \'e $\lim_{h\to 0}q(h)h=0$. Seja $K=\abs{f'(a)}+1$. Como $q(h)\to f'(a)$, existe $\delta_1>0$ tal que $0<\abs{h}<\delta_1$ implica $\abs{q(h)-f'(a)}<1$, e portanto
\[
\abs{q(h)}\le \abs{q(h)-f'(a)}+\abs{f'(a)}<1+\abs{f'(a)}=K.
\]
Dado $\eps>0$, escolha $\delta=\min(\delta_1,\eps/K)$. Se $0<\abs{h}<\delta$, ent\~ao
\[
\abs{f(a+h)-f(a)}=\abs{q(h)}\,\abs{h}<K\cdot\frac{\eps}{K}=\eps.
\]
Logo $\lim_{h\to 0}f(a+h)=f(a)$.
\end{proof}

\section{M\'odulo}

\textbf{Enunciado.} $\abs{x}$ n\~ao \'e deriv\'avel em $0$.

\begin{proof}
Para $h\neq 0$,
\[
\frac{\abs{0+h}-\abs{0}}{h}=\frac{\abs{h}}{h}.
\]
Se $h>0$, o quociente vale $1$. Se $h<0$, vale $-1$. Suponha que o limite bilateral exista e valha $L$. Dado $\eps>0$, existe $\delta>0$ tal que $0<\abs{h}<\delta$ implica $\abs{\abs{h}/h-L}<\eps$. Em particular, para $0<h<\delta$ tem-se $\abs{1-L}<\eps$. Isso vale para todo $\eps>0$. Se fosse $L\neq 1$, o n\'umero $\eps=\abs{1-L}/2>0$ daria $\abs{1-L}<\abs{1-L}/2$, o que \'e falso. Logo $L=1$. Pelo mesmo $\delta$ no lado esquerdo, $-\delta<h<0$ implica $\abs{-1-L}<\eps$ para todo $\eps>0$, e o mesmo argumento for\c{c}a $L=-1$. N\~ao se pode ter $L=1$ e $L=-1$. A derivada em $0$ n\~ao existe.
\end{proof}

\end{document}
